% =====================================================================
%  theory_appendix_additions.tex
%
%  Rigorous versions of the three deferred theoretical results for PAGE.
%  Ready to be pasted into the appendix of main.tex, AFTER the existing
%  Appendices \ref{app:proof-scaling} and \ref{app:sufficient} (they
%  reference labels defined there: eq:scaling, eq:scaling-app, eq:rho,
%  eq:D, eq:D-jaccard, eq:dilution, eq:bridge, eq:bridge-final, lem:topk,
%  and assumptions A1--A4).
%
%  This file only ADDS. It does not redefine \lemma (already in main.tex).
%  It declares the extra theorem-like environments it needs; if main.tex
%  later defines any of these, delete the duplicate \newtheorem line here.
% =====================================================================

% , - environments (main.tex already has \newtheorem{lemma}{Lemma}) , -
\newtheorem{theorem}{Theorem}
\newtheorem{proposition}{Proposition}
\newtheorem{corollary}{Corollary}
\theoremstyle{definition}
\newtheorem{assumption}{Assumption}
\theoremstyle{remark}
\newtheorem{remark}{Remark}

\section{Sharpened deferred proofs}
\label{app:sharpened}

This appendix upgrades three results that Appendices~\ref{app:proof-scaling}
and~\ref{app:sufficient} left as sketches: (i) the leading constant and a
two-sided bound for the scaling formula~\eqref{eq:scaling}; (ii) the
necessary (converse) direction of the $D$-to-dilution bridge; and (iii) the
status of the useful-set identification (A4). Each subsection ends with an
explicit ``What this does and does not establish'' paragraph so the paper's
claims can be scoped exactly. Throughout we keep the notation of
Appendices~\ref{app:proof-scaling}--\ref{app:sufficient}: $\alpha_{\mathcal{R}}$
is the pre-eviction signal mass, $\gamma$ the noise-retention ratio,
$\rho_{\mathcal{R}},\rho_{\mathcal{N}}$ the retained mass fractions, $\delta_t$
Bui's dilution~\eqref{eq:dilution}, $D$ the head-agreement drop~\eqref{eq:D},
$\sigma_z$ the pre-softmax logit-noise scale, $m$ the logit margin, $H,L,T$ the
head, layer and length counts, and $\hat A(\cdot)$ the strictly increasing,
$[0,1]$-valued, unit-saturating accuracy surrogate.

% =====================================================================
\subsection{Item 1: the exact leading constant and a two-sided scaling bound}
\label{app:tight-scaling}
% =====================================================================

The sketch in Appendix~\ref{app:proof-scaling} obtains
$\rho_{\mathrm{KV}} \le K\,(1-A_{\mathrm{full}})\,p_{\mathrm{recov}} + K/4$
by a union bound over the $K$-point budget grid $\mathcal{B}=\{b_1,\dots,b_K\}$
followed by a covariance residual. We show that the grid factor $K$ is an
artifact of an unnecessary union step, that the correct leading constant is
exactly $1$, and that a distribution-free \emph{two-sided} bound holds with
explicit endpoints; the only genuinely unresolved quantity is a covariance,
for which we give matching, attainable bounds.

\paragraph{Events.} We work inside the surrogate of
Appendix~\ref{app:proof-scaling}, where the per-input improvement event at a
single budget $b$ factorises (their displayed identity) into a headroom
indicator and a recovery indicator. Fix the prompt distribution and, for a
random prompt $x$, define the two Bernoulli random variables
\begin{equation}
\label{eq:AB-def}
A(x) := \mathbf{1}\bigl[\hat A(\mathrm{SNR}(b_{\max},x)) < 1\bigr]
\quad\text{(headroom)},
\qquad
B(x) := \mathbf{1}\Bigl[\exists\, b\in\mathcal{B}:\ \mathrm{SNR}_{\mathrm{post}}(b,x) > \mathrm{SNR}_{\mathrm{pre}}(x)\Bigr]
\quad\text{(recovery)} .
\end{equation}
The operational recovery rate is $\rho_{\mathrm{KV}} = \Pr_x[A(b^*(x),x) > A(b_{\max},x)]$
with $b^*(x)=\arg\max_b A(b,x)$ (Eq.~\eqref{eq:rho}). Because $b^*$ is the
argmax over the grid, ``improvement at $b^*$'' is exactly ``improvement at some
$b\in\mathcal{B}$,'' and the single-budget surrogate factorisation gives
\begin{equation}
\label{eq:AandB}
\bigl\{A(b^*(x),x) > A(b_{\max},x)\bigr\}
\;=\; \bigcup_{b\in\mathcal{B}}\bigl(A(x)\cap B_b(x)\bigr)
\;=\; A(x)\cap\Bigl(\bigcup_{b\in\mathcal{B}}B_b(x)\Bigr)
\;=\; A(x)\cap B(x),
\end{equation}
where $B_b(x)=\{\mathrm{SNR}_{\mathrm{post}}(b,x)>\mathrm{SNR}_{\mathrm{pre}}(x)\}$
and $B=\bigcup_b B_b$. The headroom event does not depend on $b$, so it pulls
out of the union: the union over budgets lives \emph{inside} $B$. We set
$p_{\mathrm{recov}} := \Pr_x[B]$ (recovery at the best budget) and note
$\Pr_x[A] = 1-A_{\mathrm{full}}$ by definition of the surrogate headroom.

\begin{theorem}[Exact leading constant; two-sided scaling bound]
\label{thm:tight-scaling}
With the definitions in~\eqref{eq:AB-def}--\eqref{eq:AandB} and
$p_{\mathrm{recov}}:=\Pr_x[B]$, the surrogate per-input recovery rate satisfies
the \emph{exact} identity
\begin{equation}
\label{eq:exact-scaling}
\rho_{\mathrm{KV}}
\;=\; (1-A_{\mathrm{full}})\,p_{\mathrm{recov}} \;+\; \mathrm{Cov}_x\!\bigl(A,B\bigr).
\end{equation}
The content here is not the identity~\eqref{eq:exact-scaling}, which is merely
the definition of covariance for two Bernoullis, but the observation that once
$p_{\mathrm{recov}}$ is defined as the \emph{union} probability
$\Pr_x[\bigcup_b B_b]$ (the quantity $b^*$ actually selects) rather than the
single-budget $\max_b\Pr_x[B_b]$, the budget-grid factor $K$ of
Appendix~\ref{app:proof-scaling} is absorbed into $p_{\mathrm{recov}}$ rather
than multiplying it; $K$ reappears inside $p_{\mathrm{recov}}$ whenever the
recovery budgets are near-disjoint (Corollary~\ref{cor:grid-tight}(ii)). The
covariance is bounded, by
\begin{equation}
\label{eq:cov-frechet}
0 \;\le\; \mathrm{Cov}_x(A,B) \;\le\; \min\{1-A_{\mathrm{full}},\,p_{\mathrm{recov}}\} - (1-A_{\mathrm{full}})\,p_{\mathrm{recov}},
\end{equation}
where the lower bound requires the association hypothesis of
Lemma~\ref{lem:assoc} below, and the upper bound (Fr\'echet) is
unconditional. Equivalently,
\begin{equation}
\label{eq:two-sided-scaling}
\boxed{\;
(1-A_{\mathrm{full}})\,p_{\mathrm{recov}}
\;\le\;
\rho_{\mathrm{KV}}
\;\le\;
\min\{\,1-A_{\mathrm{full}},\ p_{\mathrm{recov}}\,\}.
\;}
\end{equation}
The lower bound is attained iff $A\perp B$; the upper bound is attained iff
$A$ and $B$ are Fr\'echet--Hoeffding comonotone (one event a.s.\ contains the
other).
\end{theorem}

\begin{proof}
Identity~\eqref{eq:exact-scaling} is the definition of covariance applied to
the two Bernoulli variables in~\eqref{eq:AB-def}, using
$\rho_{\mathrm{KV}}=\Pr_x[A\cap B]=\mathbb{E}_x[AB]$ from~\eqref{eq:AandB} and
$\mathbb{E}_x[A]=1-A_{\mathrm{full}}$, $\mathbb{E}_x[B]=p_{\mathrm{recov}}$.
For~\eqref{eq:cov-frechet}, the upper endpoint is the Fr\'echet--Hoeffding
inequality $\Pr[A\cap B]\le\min\{\Pr[A],\Pr[B]\}$, valid for any two events;
subtracting the product gives the stated covariance bound and the right side
of~\eqref{eq:two-sided-scaling}. The lower endpoint $\mathrm{Cov}_x(A,B)\ge0$
is Lemma~\ref{lem:assoc}. Attainment: independence gives equality in the lower
bound by construction; comonotonicity ($A\subseteq B$ or $B\subseteq A$ a.s.)
gives $\Pr[A\cap B]=\min\{\Pr[A],\Pr[B]\}$, the upper endpoint.
\end{proof}

\begin{lemma}[Association of headroom and recovery]
\label{lem:assoc}
Suppose the noise-isotropy assumption (A1) of
Appendix~\ref{app:proof-scaling} holds, so that both $A$ and $B$ are
non-increasing functions of the single scalar driver $\alpha_{\mathcal{R}}(x)$
(equivalently, $\mathbb{E}[A\mid\alpha_{\mathcal{R}}]$ and
$\mathbb{E}[B\mid\alpha_{\mathcal{R}}]$ are non-increasing in
$\alpha_{\mathcal{R}}$ and monotone in the same direction). Then
$\mathrm{Cov}_x(A,B)\ge0$.
\end{lemma}

\begin{proof}
If $A=f(\alpha_{\mathcal{R}})$ and $B=g(\alpha_{\mathcal{R}})$ with $f,g$ both
non-increasing, Chebyshev's association (sum) inequality gives
$\mathbb{E}[fg]\ge\mathbb{E}[f]\mathbb{E}[g]$, i.e.\ $\mathrm{Cov}\ge0$. If $A$
depends on an auxiliary driver $\zeta$ independent of $\alpha_{\mathcal{R}}$
(e.g.\ a parametric-knowledge covariate), write
$\mathrm{Cov}(A,B)=\mathrm{Cov}(\mathbb{E}[A\mid\alpha_{\mathcal{R}}],B)$ by the
tower rule and $B\!=\!g(\alpha_{\mathcal{R}})$; the same monotone-in-a-common-scalar
argument applies to the conditional mean, so the sign is preserved.
\end{proof}

The monotonicity of $A$ in $\alpha_{\mathcal{R}}$ is immediate (larger signal
mass $\Rightarrow$ larger $\mathrm{SNR}_{\mathrm{pre}}$ $\Rightarrow$ smaller
headroom). The monotonicity of $B$ follows from Eq.~\eqref{eq:scaling-app}:
the attainable SNR gain from removing noise mass shrinks as
$\alpha_{\mathcal{R}}\to1$ (less noise to remove), so recovery is less likely
at larger $\alpha_{\mathcal{R}}$.

\paragraph{Why the grid factor $K$ was spurious (tight vs.\ loose regime).}
The apparent $K$ in Appendix~\ref{app:proof-scaling} comes from the step
$\Pr_x[\bigcup_b B_b]\le\sum_b\Pr_x[B_b]\le K\max_b\Pr_x[B_b]$ followed by
identifying $\max_b\Pr_x[B_b]$ with $p_{\mathrm{recov}}$. Defining
$p_{\mathrm{recov}}$ as the union probability $\Pr_x[\bigcup_b B_b]$
instead, which is the quantity the operational $b^*(x)$ actually selects, removes
the union step entirely. The overcount incurred by the discarded step is,
exactly,
\begin{equation}
\label{eq:union-overcount}
\sum_{b\in\mathcal{B}}\Pr_x[B_b] \;-\; \Pr_x\Bigl[\textstyle\bigcup_b B_b\Bigr]
\;=\; \sum_{j\ge2}(-1)^j\!\!\!\sum_{|S|=j, S\subseteq\mathcal{B}}\!\!\Pr_x\Bigl[\textstyle\bigcap_{b\in S}B_b\Bigr]
\;\ge\;0,
\end{equation}
the inclusion--exclusion tail. Its size is what makes the union bound tight or
loose:
\begin{corollary}[Tightness characterisation of the budget-grid bound]
\label{cor:grid-tight}
Let $p^\sharp:=\max_b\Pr_x[B_b]$ be the single-budget recovery probability.
\begin{enumerate}
\item[\emph{(i)}] \emph{(Monotone gain profile $\Rightarrow$ no $K$ factor.)}
If the SNR gain $g(b,x):=\rho_{\mathcal{R}}(b,x)^2/\gamma(b,x)$ is monotone in
$b$ for a.e.\ $x$ (as it is for a calibrated scorer, since keeping more tokens
raises both $\rho_{\mathcal{R}}\to1$ and $\gamma\to1$ so $g\downarrow1$), then
the events $\{B_b\}$ are nested and
$p_{\mathrm{recov}}=\Pr_x[\bigcup_b B_b]=p^\sharp$. The genuine leading factor
is $1$; the union bound $K\,p^\sharp$ overcounts by $(K-1)p^\sharp$.
\item[\emph{(ii)}] \emph{(Disjoint gain profile $\Rightarrow$ $K$ is real.)}
If instead the $\{B_b\}$ are pairwise disjoint (each input recovers at exactly
one budget) and rare ($Kp^\sharp\ll1$), then
$p_{\mathrm{recov}}=\sum_b\Pr_x[B_b]$ is genuinely $\Theta(K)$ times larger
than $p^\sharp$; recovery at the best budget really is $K$-fold more frequent
than at a fixed budget, and the factor $K$ is not slack but signal.
\end{enumerate}
In both cases Theorem~\ref{thm:tight-scaling} holds with constant $1$ once
$p_{\mathrm{recov}}$ is the union probability; the two regimes differ only in
how much larger $p_{\mathrm{recov}}$ is than the single-budget $p^\sharp$.
\end{corollary}

\begin{proof}
(i) If $g(\cdot,x)$ is monotone then $\{b:g(b,x)>1\}$ is an interval anchored at
the aggressive end of the grid, so $B_{b}(x)\subseteq B_{b'}(x)$ whenever
$b'$ is more aggressive than $b$; the $B_b$ are nested and their union equals
the loosest single event, of probability $p^\sharp$. (ii) For disjoint events
$\Pr[\bigcup_b B_b]=\sum_b\Pr[B_b]$ exactly; rareness makes this
$\approx Kp^\sharp$.
\end{proof}

Because our scorers are calibrated (the top-$k$ retention curve is monotone in
$b$), the operative regime is Corollary~\ref{cor:grid-tight}(i): the honest
constant is $1$, and the $K$ and $K/4$ terms of
Appendix~\ref{app:proof-scaling} should be read as a loose fallback that is
never binding on our grid.

\paragraph{What this does and does not establish.}
\emph{Establishes:} inside the deterministic-SNR surrogate, the scaling law
holds as the \emph{exact} identity~\eqref{eq:exact-scaling} with leading
constant precisely $1$ (no budget-grid inflation), and the two-sided
bound~\eqref{eq:two-sided-scaling} with fully explicit, distribution-free
endpoints; the only free quantity is the covariance
$\mathrm{Cov}_x(A,B)\in[0,\ \min\{\Pr A,\Pr B\}-\Pr A\Pr B]$, and both endpoints
are attainable, so no tighter distribution-free constant exists.
\emph{Does not establish:} (a) the surrogate itself, that accuracy is a
deterministic (or fixed-conditional) function of head SNR, is assumed, not
proven, so the constant ``$1$'' is exact only relative to that surrogate; (b)
the covariance is pinned only to an interval, not a value, because it depends on
the joint law of headroom and recovery, which is task- and model-specific;
Lemma~\ref{lem:assoc} fixes its \emph{sign} (non-negative) under single-driver
monotonicity but not its magnitude; (c) when headroom is driven partly by an
$\alpha_{\mathcal{R}}$-independent factor (parametric-knowledge failures, e.g.\
the QA cells in Table~\ref{tab:partition}), recovery collapses
($p_{\mathrm{recov}}\!\to\!0$) so both bounds pinch $\rho_{\mathrm{KV}}\!\to\!0$
despite headroom, consistent with the observed QA exceptions, but a prediction
of the sign structure, not a proof about those tasks. The empirically observed
ratio band $[0.25,0.55]$ is thus an estimate of
$p_{\mathrm{recov}}+\mathrm{Cov}/(1-A_{\mathrm{full}})$, not of a universal
constant.

% =====================================================================
\subsection{Item 2: the necessary direction (capacity-bound $\Rightarrow$ small $D$)}
\label{app:necessary}
% =====================================================================

Appendix~\ref{app:sufficient} proves the sufficient direction (large $D$
forces large $\delta_t$) and states that the converse ``capacity-bound
$\Rightarrow$ small $D$'' needs ``a margin condition we have not verified.''
We state that condition precisely and prove the converse under it, in both a
deterministic and a probabilistic (contrapositive) form, with explicit margin
dependence.

\paragraph{The margin condition, made precise.} Recall
$z^{(\ell)}_{h,i}=(q^{(\ell)}_h)^\top k_i + \eta^{(\ell)}_{h,i}$ with
$\sigma_z$-sub-Gaussian noise (A2), and the logit margin $m=m_+-m_-$ (A1).
Define the \emph{margin-to-noise ratio} $\kappa := m/\sigma_z$.

\begin{assumption}[$\kappa$-separation of the capacity-bound relevant set]
\label{ass:margin}
For an input in the capacity-bound regime $\mathcal{C}$, the deterministic
logit gap between every common relevant position and every noise position at
the late layer bin exceeds
\begin{equation}
\label{eq:margin-cond}
m \;\ge\; 2\,\sigma_z\sqrt{\log\!\bigl(2 H^2 L T^2/\beta\bigr)},
\qquad\text{i.e.}\qquad
\kappa \;\ge\; 2\sqrt{\log\!\bigl(2 H^2 L T^2/\beta\bigr)} .
\end{equation}
\end{assumption}
This is exactly the threshold of Lemma~\ref{lem:topk}; naming it as an
assumption isolates the one quantity the converse needs. It is the precise form
of the informal ``separation between the relevant-token logit margin and the
noise scale.''

\begin{theorem}[Capacity-bound $\Rightarrow$ small head-agreement drop]
\label{thm:converse}
Adopt the surrogate (A1)--(A3) of Appendix~\ref{app:sufficient} and
Assumption~\ref{ass:margin}. Suppose the input is capacity-bound in the sense
that the late-layer heads share a common relevant set: there is a set
$\mathcal{R}_\star$ with $\mathcal{R}^{(\ell)}_h=\mathcal{R}_\star$ for every
head $h$ and every layer $\ell$ in the late bin, and top-$k$ is taken at
$k=|\mathcal{R}_\star|$. Then, over the pre-softmax noise,
\begin{equation}
\label{eq:converse-main}
\Pr\bigl[\,a_{\mathrm{late}} = 1\,\bigr]\;\ge\;1-\beta,
\qquad\text{hence}\qquad
\Pr\bigl[\,D \le a_{\mathrm{early}}-1 \le 0\,\bigr]\;\ge\;1-\beta,
\end{equation}
and therefore, for \emph{every} threshold $\tau>0$,
\begin{equation}
\label{eq:converse-tail}
\Pr\bigl[\,D \ge \tau\,\bigr]\;\le\;\beta
\;=\;2 H^2 L T^2\,\exp\!\Bigl(-\tfrac{m^2}{4\sigma_z^2}\Bigr)
\;=\;2 H^2 L T^2\,e^{-\kappa^2/4}.
\end{equation}
\end{theorem}

\begin{proof}
By Lemma~\ref{lem:topk}, under~\eqref{eq:margin-cond} the top-$k$ set of every
late-bin head $(\ell,h)$ equals its private set $\mathcal{R}^{(\ell)}_h$
jointly with probability $\ge1-\beta$; call this event $E$. On $E$, since all
late-bin heads share $\mathcal{R}_\star$, every head's top-$k$ set is exactly
$\mathcal{R}_\star$, so every pairwise top-$k$ Jaccard equals $1$ at every
late-bin layer, whence $a_\ell=1$ for each late-bin layer and
$a_{\mathrm{late}}=1$. This proves the first claim of~\eqref{eq:converse-main}.
Consequently $D=a_{\mathrm{early}}-a_{\mathrm{late}}=a_{\mathrm{early}}-1\le0$
on $E$ because $a_{\mathrm{early}}\le1$ (Jaccard is bounded by $1$), giving the
second claim. Finally, $\{D\ge\tau\}$ for $\tau>0$ is disjoint from $E$ (on $E$,
$D\le0$), so $\Pr[D\ge\tau]\le\Pr[E^{\mathrm c}]\le\beta$; substituting the
value of $\beta$ from Lemma~\ref{lem:topk} yields~\eqref{eq:converse-tail}.
\end{proof}

Equation~\eqref{eq:converse-tail} is the explicit margin dependence: the
probability that a capacity-bound input produces a drop as large as any fixed
$\tau$ decays as $e^{-\kappa^2/4}$ in the margin-to-noise ratio. The bound is
only informative when $\beta = 2H^2 L T^2 e^{-\kappa^2/4} < 1$, which for
realistic $(H, L, T)$ requires a large margin-to-noise ratio
($\kappa \gtrsim 10$ at $H{=}16, L{=}28, T{=}4096$); this is a genuine
knife-edge (Assumption~\ref{ass:margin} is close to the informal margin cap
$m \le \log T$ of Appendix~\ref{app:proof-scaling}), and outside it the
guarantee is vacuous. Consistent with this, the measured $D$ on NIAH-MK3 is
\emph{small} rather than reliably non-positive: it is mildly positive in most
cells of Table~\ref{tab:drops} ($+0.04$ to $+0.06$) and negative only on
Qwen2.5-3B. The theorem therefore predicts $\mathbb{E}[D] \approx 0$ at an
$O(\beta)$ floor, with sign indeterminate; the small positive values observed
mean real margins sit below the high-$\kappa$ regime, so the theorem constrains
the \emph{ordering} (NIAH-MK3 smallest) more than the sign.

\begin{corollary}[Probabilistic converse / contrapositive]
\label{cor:contrapositive}
Let $\mathcal{C}_\star$ denote the capacity-bound event of
Theorem~\ref{thm:converse} (late-layer heads share a common relevant set).
Under Assumption~\ref{ass:margin}, for any $\tau>0$
\begin{equation}
\label{eq:contrapositive}
\Pr\bigl[\,\mathcal{C}_\star \ \wedge\ D\ge\tau\,\bigr]\;\le\;\beta
\;=\;2H^2 L T^2 e^{-\kappa^2/4}.
\end{equation}
Equivalently, if $\Pr[D\ge\tau]>0$ then
$\Pr[\mathcal{C}_\star \mid D\ge\tau]\le \beta/\Pr[D\ge\tau]$: observing a drop
$D\ge\tau$ implies the input is \emph{not} capacity-bound, except on a
noise-tail event of probability $\le\beta$. This is the a-priori guarantee the
gate needs: firing on $D\ge\tau$ mis-fires on a genuinely capacity-bound input
only with probability $\le\beta$.
\end{corollary}

\begin{proof}
On $\mathcal{C}_\star$, Theorem~\ref{thm:converse} gives $D\le0<\tau$ except on
$E^{\mathrm c}$; hence $\mathcal{C}_\star\cap\{D\ge\tau\}\subseteq E^{\mathrm c}$
and $\Pr[\mathcal{C}_\star\cap\{D\ge\tau\}]\le\beta$. The conditional statement
is Bayes' rule.
\end{proof}

\paragraph{What this does and does not establish.}
\emph{Establishes:} under the explicit $\kappa$-separation
Assumption~\ref{ass:margin}, the converse holds cleanly in the
surrogate, capacity-bound inputs have $D\le0$ with probability $\ge1-\beta$
(Theorem~\ref{thm:converse}), with the failure probability decaying as
$e^{-\kappa^2/4}$; and the operational contrapositive
(Corollary~\ref{cor:contrapositive}) bounds the gate's capacity-bound
mis-fire probability by the same $\beta$. This upgrades the informal ``margin
condition we have not verified'' to a named assumption with a proof and a
rate. \emph{Does not establish:} (a) that \emph{real} capacity-bound tasks
satisfy $\mathcal{C}_\star$, i.e.\ that late-layer heads literally share one
common relevant set rather than merely agreeing on an $O(1)$ neighbourhood of
the needle, top-$k$ Jaccard cannot distinguish these, and the difference is
invisible to $D$ but would weaken the deterministic $a_{\mathrm{late}}=1$ step
to $a_{\mathrm{late}}=1-o(1)$ (the conclusion $\Pr[D\ge\tau]\le\beta+o(1)$
survives, but $\mathcal{C}_\star$ is then an idealisation); (b) the numerical
value of $\kappa$ on real attention weights, Assumption~\ref{ass:margin} is a
hypothesis about the network, tested here only through its consequence (small
$D$ on NIAH-MK3 in every cell of Table~\ref{tab:drops}), not measured directly;
(c) a bound on $a_{\mathrm{early}}$ from below, so the theorem controls the sign
and the upper tail of $D$ but not how negative $D$ becomes. The converse is
thus proved \emph{conditionally on the margin assumption and the common-set
model of capacity-bound}; both are now explicit rather than hidden.

% =====================================================================
\subsection{Item 3: the A4 useful-set identification}
\label{app:a4}
% =====================================================================

Assumption (A4) of Appendix~\ref{app:sufficient} identifies Bui's useful set
$U_t$ (a residual-stream, decoding-step object) with the head-common signal
set $\mathcal{I}^{(\ell^*)}=\bigcap_h\mathcal{R}^{(\ell^*)}_h$ (a single-layer
attention object), and the text flags it as ``a modelling choice, not a derived
equality.'' We show that a clean exact equality cannot be proven, but a
\emph{quantitative approximate} identification can: under a named
consensus--usefulness assumption plus $\varepsilon$-closeness of heads, the two
sets induce dilution values that agree up to $O(\varepsilon)$, with an explicit
constant. We therefore replace the ``derived equality'' language with a proved
approximation plus a named residual assumption.

\paragraph{Set-up.} At the divergence layer $\ell^*$ write the head-private
top-$k$ sets $\mathcal{R}_h:=\mathcal{R}^{(\ell^*)}_h$, their intersection
$\mathcal{I}:=\bigcap_h\mathcal{R}_h$ and union
$\mathcal{U}:=\bigcup_h\mathcal{R}_h$. For a set $\mathcal{S}$ define its
induced dilution $\delta(\mathcal{S}):=1-\sum_{i\in\mathcal{S}}\alpha_{t,i}$, so
Bui's is $\delta^{(U)}:=\delta(U_t)$ and ours is
$\delta^{(I)}:=\delta(\mathcal{I})$. Let $\bar\alpha_{\max}:=\max_i\alpha_{t,i}$.

\begin{assumption}[$\varepsilon$-closeness of heads]
\label{ass:eps-close}
The late-layer heads are $\varepsilon$-close in their top-$k$ sets:
$\max_{h,h'}\,|\mathcal{R}_h\,\triangle\,\mathcal{R}_{h'}|\le\varepsilon k$,
equivalently pairwise Jaccard $\ge (1-\varepsilon/2)/(1+\varepsilon/2)=1-O(\varepsilon)$.
\end{assumption}

\begin{assumption}[Consensus--usefulness, CUA]
\label{ass:cua}
The useful set is sandwiched between the head-consensus and head-union sets:
$\mathcal{I}\subseteq U_t\subseteq\mathcal{U}$. The right inclusion
($U_t\subseteq\mathcal{U}$) holds whenever routing to the residual stream is
through attention only, so a position no head attends to cannot affect the
next-token margin. The left inclusion ($\mathcal{I}\subseteq U_t$) is the
substantive content: a position every head attends to is margin-preserving.
\end{assumption}

\begin{proposition}[Quantitative approximate identification]
\label{prop:a4}
Under Assumptions~\ref{ass:eps-close}--\ref{ass:cua},
\begin{equation}
\label{eq:a4-bound}
\bigl|\,\delta^{(U)}_t-\delta^{(I)}_t\,\bigr|
\;\le\;\sum_{i\in\mathcal{U}\setminus\mathcal{I}}\alpha_{t,i}
\;\le\; H(H-1)\,\varepsilon\,k\,\bar\alpha_{\max}
\;=\; O\!\bigl(H^2\varepsilon\bigr),
\end{equation}
where the last step uses $\bar\alpha_{\max}=\Theta(1/k)$ in the high-SNR
softmax. In particular, as $\varepsilon\to0$ (heads perfectly agree) the two
dilution values coincide: $\delta^{(U)}_t\to\delta^{(I)}_t$.
\end{proposition}

\begin{proof}
By Assumption~\ref{ass:cua}, $U_t\triangle\mathcal{I}\subseteq\mathcal{U}\setminus\mathcal{I}$,
because $\mathcal{I}\subseteq U_t\subseteq\mathcal{U}$ forces every element of
$U_t\triangle\mathcal{I}=U_t\setminus\mathcal{I}$ to lie in
$\mathcal{U}\setminus\mathcal{I}$. Since $\delta(\mathcal{S})$ is a signed sum of
attention masses over $\mathcal{S}$,
$|\delta^{(U)}_t-\delta^{(I)}_t|
=\bigl|\sum_{i\in U_t}\alpha_{t,i}-\sum_{i\in\mathcal{I}}\alpha_{t,i}\bigr|
\le\sum_{i\in U_t\triangle\mathcal{I}}\alpha_{t,i}
\le\sum_{i\in\mathcal{U}\setminus\mathcal{I}}\alpha_{t,i}$,
the first inequality of~\eqref{eq:a4-bound}. For the cardinality bound, for each
head $h$,
$|\mathcal{R}_h\setminus\mathcal{I}|\le\sum_{h'\ne h}|\mathcal{R}_h\setminus\mathcal{R}_{h'}|
\le\sum_{h'\ne h}|\mathcal{R}_h\triangle\mathcal{R}_{h'}|\le(H-1)\varepsilon k$
by Assumption~\ref{ass:eps-close}. Since
$\mathcal{U}\setminus\mathcal{I}=\bigcup_h(\mathcal{R}_h\setminus\mathcal{I})$,
$|\mathcal{U}\setminus\mathcal{I}|\le\sum_h|\mathcal{R}_h\setminus\mathcal{I}|\le H(H-1)\varepsilon k$.
Bounding each mass by $\bar\alpha_{\max}$ gives the middle inequality;
$\bar\alpha_{\max}=\Theta(1/k)$ gives the $O(H^2\varepsilon)$ order.
\end{proof}

\begin{remark}[When the identification fails]
\label{rem:a4-fail}
The left inclusion of Assumption~\ref{ass:cua} ($\mathcal{I}\subseteq U_t$) can
fail, and this is exactly the case the equality cannot cover: if a
\emph{single} specialised head carries the decisive value for the next-token
margin, then a position in $\mathcal{R}_{h_0}\setminus\mathcal{I}$ (attended by
that one head, absent from the consensus) can be margin-critical, so
$U_t\not\subseteq\mathcal{U}\setminus(\mathcal{R}_{h_0}\setminus\mathcal{I})$ and
$U_t\ne\mathcal{I}$ even at $\varepsilon=0$ of the \emph{other} heads. This is
the single-decisive-head regime, precisely a capacity-bound, single-needle
retrieval situation, where one retrieval head dominates. The approximation of
Proposition~\ref{prop:a4} is therefore reliable in the dilution-prone regime
(many heads, diffuse consensus, no single decisive head) and degrades exactly
where the paper already declines to apply the bridge (capacity-bound).
\end{remark}

\paragraph{What this does and does not establish.}
\emph{Establishes:} the identification of $U_t$ with $\mathcal{I}^{(\ell^*)}$ is
not an exact derived equality, but under the named
Assumptions~\ref{ass:eps-close}--\ref{ass:cua} the two sets induce
\emph{dilution values} that agree up to $O(H^2\varepsilon)$ in attention mass
(Proposition~\ref{prop:a4}), with an explicit constant $H(H-1)k\bar\alpha_{\max}$;
so wherever heads are $\varepsilon$-close, using $\mathcal{I}^{(\ell^*)}$ in
place of $U_t$ perturbs Bui's $\delta_t$ by a controlled $O(\varepsilon)$ term
rather than by an uncontrolled amount. \emph{Does not establish:} (a) the
substantive left inclusion $\mathcal{I}\subseteq U_t$ (CUA); it is named as
Assumption~\ref{ass:cua}, justified by the additive-routing picture, and shown
in Remark~\ref{rem:a4-fail} to fail precisely in the single-decisive-head
(capacity-bound) regime, so it is an assumption about the model's read-out, not
a theorem; (b) any control of the $H^2$ prefactor, which is a worst-case union
count and is loose when residuals overlap across heads (the same $O(H)$-slack
issue as Step~3 of Appendix~\ref{app:sufficient}); (c) a bound when
$\varepsilon$ is $\Theta(1)$ (genuinely divergent heads, the dilution-prone late
layers themselves), there the approximation is vacuous and $\mathcal{I}$ and
$U_t$ are related only through the CUA sandwich, not quantitatively. The honest
scoping is thus: \textbf{A4 is a named modelling assumption (CUA), under which
we prove an $O(H^2\varepsilon)$ approximate identification of the dilution
values}, not a derived equality; the word ``derived equality'' in
Appendix~\ref{app:sufficient} should be read as this approximation.
